% Part: proof-theory
% Chapter: normalization
% Section: reductions

\documentclass[../../../include/open-logic-section]{subfiles}

\begin{document}

\olfileid{pt}{nor}{red}

\olsection{ઘટાડા-રૂપાંતરો}

\Log{N1i}ની
!!a{proof}માં કટની
લંબાઈ~$1$ હોય,
તો તેનો ખંડ
માત્ર એક
!!{formula} આવૃત્તિનો
બને છે. એ
!!{formula} એક
!!a{introduction}
અનુમાનનું નિષ્કર્ષ
અને !!a{elimination}
અનુમાનનું મુખ્ય
આધારવિધાન બંને
છે. સામાન્યીકરણની
સાબિતીમાં, આવા
!!{introduction}/!!{elimination}
યુગ્મમાં પૂરી થતી
ઉપ-!!{proof}ના સ્થાને,
તે યુગ્મ કાઢી
નાખેલી નવી
!!{proof} મૂકીને
કટ ``ઘટાડીએ'' છીએ.
આ બદલાવથી નવા
કટ ઊભા થઈ
શકે છે. જોકે
નવા કટની કક્ષા
ઘટાડાતા કટની
કક્ષા કરતાં ઓછી
હોય છે. તેથી
નવી !!{proof}ની
કટ-લંબાઈ ઓછી
હોય છે (કારણ
કે લંબાઈ~$1$નો
મહત્તમ કટ દૂર
થયો), અથવા
!!{proof}ની કટ-કક્ષા
ઘટે છે (જો એ
જ એકમાત્ર
મહત્તમ કટ હોય).

દાખલા તરીકે,
સંબંધિત અનુમાનો
\Intro{\lif} અને
\Elim{\lif} હોય,
$!A \ident !B \lif !C$
હોય અને
ઉપ-!!{proof}~$\delta_1$નો
અંત આ રીતે
થાય:
\[
\AxiomC{$\Discharge{!B}{x}$}
\RightLabel{$\delta_2$}
\DeduceC{$!C$}
\DischargeRule{\Intro{\lif}}{x}
\UnaryInfC{$!B \lif !C$}
\AxiomC{}
\RightLabel{$\delta_3$}
\DeduceC{$!B$}
\RightLabel{\Elim{\lif}}
\BinaryInfC{$!C$}
\DisplayProof
\]
આ કટ દૂર કરવા
$\delta$માંની
ઉપસાબિતી~$\delta_1$ના
સ્થાને આ મૂકીએ:
\[
\AxiomC{}
\RightLabel{$\delta_3$}
\DeduceC{$!B$}
\RightLabel{$\Subst{\delta_2}{\delta_3}{!B^x}$}
\DeduceC{$!C$}
\DisplayProof
\]
$\Subst{\delta_2}{\delta_3}{!B^x}$
એટલે~$\delta_2$માં
$x$ ચિહ્નવાળી
દરેક ધારણા~$!B$ને
આખી !!{proof}~$\delta_3$થી
બદલતાં મળતી
!!{proof}. આ
!!{proof}માં
$x$ ચિહ્નવાળી
$!B$ સ્વરૂપની
ખુલ્લી ધારણા
રહેતી નથી.
તેમાં~$x$
ચિહ્નવાળી બીજી
કોઈ ધારણા પણ
નથી, કારણ કે
હંમેશની જેમ
અલગ !!{formula}s
એક જ ચિહ્ન
સાથે ધારણા
તરીકે આવતાં
નથી એમ માનીએ
છીએ. તેથી નવી
!!{proof}ની ખુલ્લી
ધારણાઓ~$\delta_1$ની
જેટલી જ છે.
$\delta_2$માં ધારણાઓ
મુક્ત કરતું
દરેક અનુમાન
$\Subst{\delta_2}{\delta_3}{!B^x}$
માં પણ એ જ
ધારણાઓ મુક્ત
કરે છે; અને
$\delta_3$માં ધારણાઓ
મુક્ત કરતું
દરેક અનુમાન
શક્ય હોય એવી
અનેક નકલોમાં
આવે છે, જે
$\Subst{\delta_2}{\delta_3}{!B^x}$
માં અનુરૂપ
ધારણાઓ મુક્ત
કરે છે.

સંપૂર્ણપણે
$\delta_2$ અથવા
$\delta_3$ની અંદરનો
કટ, કે
$\delta$માં~$\delta_1$ની
સંપૂર્ણપણે બહારનો
કટ, યથાવત્
રહે છે: તે
જ !!{formula}s
અને એ જ
અનુમાનોમાંથી
પસાર થાય છે.
તેથી ખાસ કરીને,
$\delta$ના અનુરૂપ
કટ જેટલી જ
તેની કક્ષા અને
લંબાઈ છે.
આપણે લંબાઈ~$1$નો
મહત્તમ કટ
દૂર કર્યો છે;
એટલે કટ-લંબાઈ
અથવા કટ-કક્ષા
ઘટાડી છે (જો
$\delta$ની કટ-લંબાઈ~$1$
હતી અને આ
એકમાત્ર મહત્તમ
કટ હતો).
\sourcecorrection{OLMD-255}{મૂળ ફકરો $\delta_4$ની અંદરના કટનો ઉલ્લેખ કરે છે, પરંતુ આ અનુસરણ-પરિચય/લોપના ઉદાહરણમાં ફક્ત $\delta_2$ અને $\delta_3$ ઉપસાબિતીઓ છે.}

જોકે બે બાબતો
બની શકે છે:
\begin{enumerate}
  \item બધી
  ધારણાઓ~$!B^x$ને
  $\delta_3$થી બદલતાં
  $\delta_3$ના બધા
  કટની નકલો વધી
  જાય છે. એમાંથી
  કેટલાક મહત્તમ
  કટ હોય તો
  નવી !!{proof}ની
  કટ-લંબાઈ જૂની
  !!{proof}ની કટ-લંબાઈ
  કરતાં વધુ થઈ
  શકે. આવું
  ન થાય તેની
  ખાતરી જોઈએ.
  તેથી ઘટાડા-રૂપાંતર
  ફક્ત \emph{સૌથી
  જમણા} મહત્તમ
  કટ પર લગાવીએ:
  એટલે~$\delta$ની
  શાખાઓમાં ઘટાડાતા
  કટની જમણી બાજુ
  બીજો કોઈ
  મહત્તમ કટ ન
  હોય. $\delta_3$માં
  કટ હોય તો
  તે અહીં ઘટાડાતા
  કટની જમણી
  બાજુએ હશે;
  તેથી પસંદ કરેલો
  કટ સૌથી જમણો
  ન ગણાય. હંમેશાં
  સૌથી જમણો
  કટ પસંદ કરતાં
  !!{proof}ની
  કટ-લંબાઈ અથવા
  કટ-કક્ષા ઘટે
  છે તેની ખાતરી
  મળે છે.
  \item જો
  $\delta_2$માંની
  ધારણા~$!B^x$
  !!a{elimination}
  અનુમાનનું મુખ્ય
  આધારવિધાન હોય,
  અને~$\delta_3$નું
  છેલ્લું અનુમાન
  \Elim{\lor},
  \Elim{\lexists}
  અથવા !!a{introduction}
  હોય, તો આ
  $\delta_2$માંની આ
  ધારણાની જગ્યાએ
  $\delta_3$નું
  કલમ-જોડાણ
  કરવાથી નવો
  કટ ઊભો થાય
  છે. $\delta_2$માં
  જે માત્ર ધારણા
  હતી તે હવે
  લંબાઈ~$1$નો
  કટ (એક
  !!a{introduction}
  અનુમાનનું નિષ્કર્ષ
  અને !!a{elimination}
  અનુમાનનું મુખ્ય
  આધારવિધાન) અથવા
  કટ-ખંડની
  છેલ્લી !!{formula}
  આવૃત્તિ
  (!!a{elimination}
  અનુમાનનું
  મુખ્ય આધારવિધાન
  અને~\Elim{\lor}
  કે~\Elim{\lexists}નું
  નિષ્કર્ષ) બને
  છે. જો~$!B^x$
  એ~\Elim{\lor}
  અથવા~\Elim{\lexists}નું
  ગૌણ આધારવિધાન
  હતું, તો તે
  પહેલેથી કોઈ
  ખંડની, અને
  કદાચ કટ-ખંડની,
  પ્રથમ !!{formula}
  આવૃત્તિ હતું.
  નવી !!{proof}માં
  એ ખંડ હવે
  લાંબો થઈ શકે.
  પણ નોંધો કે
  આ નવા કટ
  અથવા જૂના
  કટ-ખંડનું
  !!{formula}~$!B$
  છે અને
  $\depth{!B} < \depth{!B \lif !C}$.
  તેથી કટ ઉમેરાયા
  હોય અથવા તેમની
  લંબાઈ વધી હોય
  તો પણ તેઓ
  મહત્તમ કટ નથી.
  આમ કટ-કક્ષા
  ઘટે છે; અથવા
  સમાન કક્ષાના
  કટ બાકી હોય
  તો કટ-લંબાઈ
  ઘટે છે.
  \sourcecorrection{OLMD-256}{મૂળમાં આ જ ચિહ્નિત ધારણા પહેલા $!B^x$ અને પછી $B^x$ લખાઈ છે; આખા ઉદાહરણમાં સૂત્રનું ચિહ્ન $!B^x$ જ છે.}
\end{enumerate}

લંબાઈ~$1$નો
કટ ઘટાડવાની
પ્રક્રિયાને
\emph{ઘટાડા-રૂપાંતર}
(અથવા
\emph{ચકરાવા-રૂપાંતર})
કહે છે.
કેટલાક નિયમ-યુગ્મો
માટેના નમૂના
\olref{tab:reduction-conversions}માં
આપ્યાં છે.
(બાકીનાં કસરત
તરીકે મૂક્યાં છે.)

\begin{table}
\begin{multline*}
\bottomAlignProof
\AxiomC{$\Discharge{!B}{x}$}
\RightLabel{$\delta_2$}
\shortDeduce
\DeduceC{$!C$}
\DischargeRule{\Intro{\lif}}{x}
\UnaryInfC{$!B \lif !C$}
\AxiomC{}
\RightLabel{$\delta_3$}
\shortDeduce
\DeduceC{$!B$}
\RightLabel{\Elim{\lif}}
\BinaryInfC{$!C$}
\DisplayProof
\quad \longrightarrow\quad
\shoveright{\bottomAlignProof
\AxiomC{}
\RightLabel{$\delta_3$}
\shortDeduce
\DeduceC{$!B$}
\RightLabel{$\Subst{\delta_2}{\delta_3}{!B^x}$}
\shortDeduce
\DeduceC{$!C$}
\DisplayProof}
\\[4ex]
\shoveleft{\bottomAlignProof
\AxiomC{}
\RightLabel{$\delta_2$}
\shortDeduce
\DeduceC{$!B$}
\AxiomC{}
\RightLabel{$\delta_3$}
\shortDeduce
\DeduceC{$!C$}
\RightLabel{\Intro{\land}}
\BinaryInfC{$!B \land !C$}
\RightLabel{\Elim{\land}[1]}
\UnaryInfC{$!B$}
\DisplayProof
\quad \longrightarrow\quad
\bottomAlignProof
\AxiomC{}
\RightLabel{$\delta_2$}
\shortDeduce
\DeduceC{$!B$}
\DisplayProof}\qquad
\shoveright{\bottomAlignProof
\AxiomC{}
\RightLabel{$\delta_2$}
\shortDeduce
\DeduceC{$!B$}
\AxiomC{}
\RightLabel{$\delta_3$}
\shortDeduce
\DeduceC{$!C$}
\RightLabel{\Intro{\land}}
\BinaryInfC{$!B \land !C$}
\RightLabel{\Elim{\land}[2]}
\UnaryInfC{$!C$}
\DisplayProof
\quad \longrightarrow\quad
\bottomAlignProof
\AxiomC{}
\RightLabel{$\delta_3$}
\shortDeduce
\DeduceC{$!C$}
\DisplayProof}
\\[4ex]
\bottomAlignProof
\AxiomC{}
\RightLabel{$\delta_2$}
\shortDeduce
\DeduceC{$!B$}
\RightLabel{\Intro{\lor}[1]}
\UnaryInfC{$!B \lor !C$}
\AxiomC{$\Discharge{!B}{x}$}
\RightLabel{$\delta_3$}
\shortDeduce
\DeduceC{$!D$}
\AxiomC{$\Discharge{!C}{y}$}
\RightLabel{$\delta_4$}
\shortDeduce
\DeduceC{$!D$}
\DischargeRule{\Elim{\lor}}{x\,y}
\TrinaryInfC{$!D$}
\DisplayProof
\quad \longrightarrow\quad
\shoveright{\bottomAlignProof
\AxiomC{}
\RightLabel{$\delta_2$}
\shortDeduce
\DeduceC{$!B$}
\RightLabel{$\Subst{\delta_3}{\delta_2}{!B^x}$}
\shortDeduce
\DeduceC{$!D$}
\DisplayProof}
\\[2ex]
\shoveleft{\bottomAlignProof
\AxiomC{}
\RightLabel{$\delta_2$}
\shortDeduce
\DeduceC{$!B(c)$}
\RightLabel{\Intro{\lforall}}
\UnaryInfC{$\lforall[x][!B(x)]$}
\RightLabel{\Elim{\lforall}}
\UnaryInfC{$!B(t)$}
\DisplayProof
\quad \longrightarrow\quad
\bottomAlignProof
\AxiomC{}
\RightLabel{$\Subst{\delta_2}{t}{c}$}
\shortDeduce
\DeduceC{$!B(t)$}
\DisplayProof}
\\[4ex]
\shoveright{\bottomAlignProof
\AxiomC{}
\RightLabel{$\delta_2$}
\shortDeduce
\DeduceC{$!B(t)$}
\RightLabel{\Intro{\lexists}}
\UnaryInfC{$\lexists[x][!B(x)]$}
\AxiomC{$\Discharge{!B(c)}{x}$}
\RightLabel{$\delta_3$}
\shortDeduce
\DeduceC{$!D$}
\DischargeRule{\Elim{\lexists}}{x}
\BinaryInfC{$!D$}
\DisplayProof
\quad \longrightarrow\quad
\bottomAlignProof
\AxiomC{}
\RightLabel{$\delta_2$}
\shortDeduce
\DeduceC{$!B(t)$}
\RightLabel{$\Subst{\Subst{\delta_3}{t}{c}}{\delta_2}{!B(t)^x}$}
\shortDeduce
\DeduceC{$!D$}
\DisplayProof}
\end{multline*}
\caption{\Log{N1i} માટેનાં કેટલાંક ઘટાડા-રૂપાંતરો.}
\ollabel{tab:reduction-conversions}
\end{table}

\begin{prob}
\Intro{\lnot} પછી
\Elim{\lnot} આવે
ત્યારેનું
ઘટાડા-રૂપાંતર
આપો.
\end{prob}

હજી એક કિસ્સો
જોવાનો બાકી છે.
કટની વ્યાખ્યામાં
લંબાઈ~$1$ હોય
અને~$!A$ એ
\FalseInt{}નું
નિષ્કર્ષ તથા
!!a{elimination}
અનુમાનનું મુખ્ય
આધારવિધાન હોય,
તે કિસ્સો પણ
આવે છે. $!A$
એ !!a{introduction}
નિયમનું નિષ્કર્ષ
હોય તે કિસ્સાની
જેમ આને પણ
સંભાળીએ. દાખલા
તરીકે, આને
\[
\bottomAlignProof
\AxiomC{}
\RightLabel{$\delta_2$}
\DeduceC{$\lfalse$}
\RightLabel{\FalseInt}
\UnaryInfC{$!B \lif !C$}
\AxiomC{}
\RightLabel{$\delta_3$}
\DeduceC{$!B$}
\RightLabel{\Elim{\lif}}
\BinaryInfC{$!C$}
\DisplayProof
\quad\text{ના સ્થાને}\quad
\bottomAlignProof
\AxiomC{}
\RightLabel{$\delta_2$}
\DeduceC{$\lfalse$}
\RightLabel{\FalseInt}
\UnaryInfC{$!C$}
\DisplayProof
\]
\sourcecorrection{OLMD-257}{મૂળના પ્રથમ અસત્ય-પ્રવેશ રૂપાંતરના બે વૃક્ષોમાં અસત્યમાંથી નવું સૂત્ર કાઢતી એકપદી કડીઓ પર નિયમનું લેબલ છૂટ્યું છે; નીચેના સમાન દાખલાની જેમ \FalseInt{} દેખાડ્યું છે.}
$!A \ident (!B \lor
!C)$ અથવા
$!A \ident \lexists[x][!B(x)]$
હોય ત્યારે થોડી
સાવચેતી જોઈએ.
દાખલા તરીકે,
જો આને
\[
\bottomAlignProof
\AxiomC{}
\RightLabel{$\delta_2$}
\DeduceC{$\lfalse$}
\RightLabel{\FalseInt}
\UnaryInfC{$!B \lor !C$}
\AxiomC{$\Discharge{!B}{x}$}
\RightLabel{$\delta_3$}
\DeduceC{$!D$}
\AxiomC{$\Discharge{!C}{y}$}
\RightLabel{$\delta_4$}
\DeduceC{$!D$}
\DischargeRule{\Elim{\lor}}{x\,y}
\TrinaryInfC{$!D$}
\DisplayProof
\quad\text{ના સ્થાને}\quad
\bottomAlignProof
\AxiomC{}
\RightLabel{$\delta_2$}
\DeduceC{$\lfalse$}
\RightLabel{\FalseInt}
\UnaryInfC{$!D$}
\DisplayProof
\]
\sourcecorrection{OLMD-258}{મૂળના વિકલ્પ-લોપ વૃક્ષમાં $!B$ અને $!C$ની જુદી ધારણાઓ બંનેને $x$ ચિહ્ન આપ્યું છે અને લોપ-નિયમ પર મુક્તિ-ચિહ્નો નથી. અગાઉના N1 નિયમવૃક્ષ પ્રમાણે બીજી ધારણા $y$ અને નિયમ $x,y$ બંનેને મુક્ત કરે છે.}
તો કટનું
!!{formula}~$!D$
ધરાવતો નવો
કટ ઊભો થઈ
શકે છે. પરંતુ
$\depth{!D} < \depth{!B \lor !C}$
હોવાની કોઈ
ખાતરી નથી!{}
તેથી અહીં પણ
\Intro{\lor}/\Elim{\lor}
ઘટાડા-રૂપાંતરની
જેમ આગળ વધવું
અને તેના
સ્થાને આ મૂકવું:
\[
\bottomAlignProof
\AxiomC{}
\RightLabel{$\delta_2$}
\DeduceC{$\lfalse$}
\RightLabel{\FalseInt}
\UnaryInfC{$!B$}
\RightLabel{$\delta_3$}
\DeduceC{$!D$}
\DisplayProof
\quad\text{અથવા}\quad
\bottomAlignProof
\AxiomC{}
\RightLabel{$\delta_2$}
\DeduceC{$\lfalse$}
\RightLabel{\FalseInt}
\UnaryInfC{$!C$}
\RightLabel{$\delta_4$}
\DeduceC{$!D$}
\DisplayProof
\]

\end{document}
